% Full proof, owner routing, and exact source identities accompany this module.
\subsection{From signing reservations to ordinary unforgeability}
\label{sec:reservation-ordinary-bridge}

The retained randomizer can be incorporated into an ordinary-signature
reduction. The simulator preserves every opened reservation and uses the
complete pair of original randomizer and public message representative to
identify repeated computations. For native ML-DSA, a shared-oracle
comparison also accounts for the dependence created by key generation.

Use the owner interface of
Proposition~\ref{prop:reservation-signature-components}. A fresh owner
generation receives a family index; exact references and imported copies
retain that index. Independently generated owners remain separate when
their public keys coincide. The family ledger records each original
message at its first private invocation, including failed publication.
It is contained in the native actual-key ledger. Stopping before a
target-family key or signer-file return therefore gives an auxiliary
fresh-message event containing the native stopped event. All routing and
stopping decisions use declared lineage and public metadata, without
testing hidden secret equality.

\paragraph{Independent hedging.}
Suppose signing factors as
\[
 \operatorname{Sign}(s,K,M;r)
 =S\bigl(s,d_{pk}(M);F_K(r,d_{pk}(M))\bigr),
\]
where the complete computation after the public descriptor $d_{pk}(M)$
is deterministic given its internal tape, and $K$ is generated
independently of the remaining owner state. The ordinary signer draws a
fresh hidden $r\in\{0,1\}^n$ per call. The reservation interface reveals
original randomizers before or after execution and retains the complete
original-message ledger. The joint PRF experiment supplies the same public
oracle access as the signature game, including quantum queries when
present; its private function calls are classical and its ideal random
function is independent of the public companion oracle.

\begin{proposition}[An ordinary-signature bridge for independent hedging]
\label{prop:independent-hedging-bridge}
For one protected family and at most $q$ distinct executed
$(r,d_{pk}(M))$ pairs, there are an ordinary wrapper-EUF adversary $B$ and
joint-PRF distinguishers $D_0,D_1$ such that
\begin{equation}
 p_{\rm res}(A)\le p_{\rm EUF}(B)
   +\operatorname{Adv}^{\rm jprf}_F(D_0)
   +\operatorname{Adv}^{\rm jprf}_F(D_1)
   +\beta_n(q),\qquad
 \beta_n(q)=\min\{1,\tbinom q2 2^{-n}\}.
 \label{eq:independent-hedging-bridge}
\end{equation}
Each primitive interface receives at most $q$ calls. The bounds include
the complete simulation, public descriptor work and owner observations.
\end{proposition}
\begin{proof}
Replace the reservation signer's hedging function by a private random
function, using independence to simulate the remaining key state. This
costs the first joint-PRF advantage. Simulator $B$ samples all visible
reservation randomizers and caches signatures by the full pair
$(r,d_{pk}(M))$. A new pair causes one ordinary signing query on its
original $M$; every first slot execution enters its original message in
the local ledger, including a failed publication. A repeated pair reuses
the cached signature.

Replace the ordinary challenger's hedge by a private random function.
Until its fresh hidden input repeats, its internal tapes are independent
uniform values and can be coupled with the tapes of the distinct visible
pairs. The $t$-th hidden draw repeats a prior input with probability at
most $(t-1)2^{-n}$, giving $\beta_n(q)$. This charge concerns the ordinary
challenger's draws. Arbitrarily many opened reservations and adaptively
selected visible collisions are represented exactly by the pair cache.
Returning the ordinary challenger to its real hedge costs the second
joint-PRF term. Its signing queries are a subset of the executed-message
ledger, so a fresh output is an ordinary forgery. Descriptor aliases need
no separate charge: factorization gives the same signature, while the
ledger retains every original message.
\end{proof}

For several protected families with a global stop before any of their
key/file disclosures, the same argument uses indexed multiuser primitive
games and the birthday term
$\min\{1,2^{-n}\sum_h\binom{q_h}{2}\}$, with no extra user factor.
Permitting other-owner disclosures instead uses the selected-family
construction below. All displayed query caps hold across the relevant
complete hybrids; an Execute-call cap is always a valid conservative
choice for $q$.

\paragraph{The shared ML-DSA oracle.}
FIPS~204 key generation expands a uniform 32-byte seed into consecutive
blocks $(\rho,\rho',K)$ of lengths $(32,64,32)$ bytes. Signing computes
\[
 \mu=H(\mathit{tr}\Vert M',64),\qquad
 c=H(K\Vert r\Vert\mu,64).
\]
The remaining signing computation uses the other key components, $\mu$
and $c$~\cite{fips204}. The pinned native backend passes its retained
32 bytes to \texttt{try\_sign\_with\_seed}; the checked implementation
uses $K$ only for this hedge in the signature computation. The full source
map records the fixed crate archive and its upstream commit.

Model SHAKE256 as a public random XOF $H$, with one independent infinite
output string per complete input and consistent replies to every requested
prefix. Other public algorithms retain the same interfaces in the
comparison and are independent of the sampled random-XOF table. For a
selected ML-DSA-65 family, let $q_*^{\rm res}$ and
$q_*^{\rm ord}$ bound all remaining calls that can address a 34-byte or
128-byte input in the two complete experiments, after rewriting every
designated target seed expansion and hedge as in the proof. Count public
adversarial queries, preparation, verification and all other-owner work.
Define
\[
 \delta_{\rm C}(v)=\min\{1,v2^{-256}\},\qquad
 \delta_{\rm Q}(v)=\min\{1,2v2^{-128}\}.
\]

\begin{theorem}[Native ML-DSA reservation bridge]
\label{thm:native-mldsa-reservation-bridge}
For the declared owner interface and a selected native ML-DSA-65 family,
let $q$ bound the distinct executed target $(r,\mu)$ pairs in every
relevant experiment. There is an ordinary hedged-ML-DSA EUF adversary $B$ with at most $q$
signing queries such that
\begin{equation}
 p_{\rm res}^{H}(A)
 \le p_{\rm EUF}^{H}(B)+\beta_{256}(q)
       +\delta_{\mathsf X}(q_*^{\rm res})
       +\delta_{\mathsf X}(q_*^{\rm ord}).
 \label{eq:native-mldsa-reservation-bridge}
\end{equation}
Here $\mathsf X=\mathsf C$ denotes the classical random-oracle model and
$\mathsf X=\mathsf Q$ the quantum random-oracle model. Reservation and
signing calls are classical in both. Non-target owners may be generated
locally and exposed; the target interface stops before key or signer-file
return and has no earlier secret-byte observation. Both sides use the
same ordinary native algorithm and public-oracle model, with their full
simulation time and oracle budgets.
\end{theorem}
\begin{proof}
Sample independent uniform $\xi,\rho,\rho',K$ of the specified lengths,
a random XOF $H_0$, and a private random function
$R:(r,\mu)\mapsto\{0,1\}^{512}$. Let
\[
 g=\xi\Vert\operatorname{byte}(6)\Vert\operatorname{byte}(5),\qquad
 \mathcal S_K=\{K\Vert r\Vert\mu:\ |r|=32,\ |\mu|=64\},
\]
where these lengths are in bytes. Define $H_*$ by replacing the first
128 output bytes at $g$ with $\rho\Vert\rho'\Vert K$, and the first
64 bytes at each $K\Vert r\Vert\mu$ with $R(r,\mu)$. Preserve every
remaining output bit from $H_0$. The two input sets have disjoint lengths
34 and 128. Conditional on $\xi$ and the keygen output, all other full XOF strings
remain independent uniform, including those in the selected slice.
Thus $H_*$ and the keygen blocks have exactly the native joint distribution.

Rewrite every target-lineage seed expansion, including reopened owners,
to use those sampled blocks, and every target hedge, including retries,
to use $R$. These are exact substitutions under $H_*$. Change every
remaining oracle call to $H_0$. In this ideal experiment the query-register
distribution is independent of the now-unused $\xi,K$. The remaining
base state may depend on $\rho,\rho',H_0$, and every non-target owner keeps
its original algorithm and shared-oracle calls.

For any ideal query, its expected mass on
$\mathcal D=\{g\}\cup\mathcal S_K$ is at most $2^{-256}$: each relevant
length sector has that hidden-point or hidden-prefix probability, and the
two sector masses sum to at most one. A classical identical-until-hit
argument gives $q_*2^{-256}$. For a quantum query, the difference between
the two XOR oracle actions has norm at most $2\sqrt{p_t(\mathcal D)}$.
Telescoping using the ideal prefix states and applying Jensen's inequality
gives expected final-state distance at most $2q_*2^{-128}$, which also
bounds the difference in a final event probability
\cite{bennett1997strengths}. This comparison includes verification.

Apply it first to the reservation experiment. In the resulting ideal
experiment $R$ is independent of the public oracle, so the pair-cache
coupling from Proposition~\ref{prop:independent-hedging-bridge} gives the
ordinary ideal experiment with loss $\beta_{256}(q)$. Apply the same
oracle comparison in the reverse direction to the complete ordinary
experiment. Its signer is now the native ordinary challenger, giving
\eqref{eq:native-mldsa-reservation-bridge}. The actual reduction $B$ uses
only that signing oracle and public interfaces; the statistical comparison
does not require it to know or program $\xi$ or $K$.
\end{proof}

The source inventory makes the query budget concrete. After the designated
rewrites, native ML-DSA-65 uses SHAKE256 input lengths 66 for secret-vector
and mask expansion, 48 for challenge-to-ball sampling, 832 for commitment
hashes, and 1952 for the public-key hash. Matrix expansion uses SHAKE128.
The protocol's fixed context is 42 bytes; its complete descriptor input
has length $155+|B|$ bytes. These target computations therefore agree
exactly under $H_*$ and $H_0$, including every rejection-loop iteration.
Arbitrary public SHAKE256 calls and non-target seed/hedge calls remain in
$q_*$. Safe-domain calls are common operations in the query hybrid. A
finite computation cutoff followed by a limit proves the same statistical
bound for unbounded safe-domain rejection loops; full running time remains
part of the ordinary-EUF term.

\paragraph{Adaptive owners and the pair signature.}
For at most $U$ fresh owner families, guess the target's creation index
uniformly, retaining that index across aliases. Embed the single-user
challenger there and generate all other owners locally, answering their
permitted disclosures. This works for adaptive creation and gives
\begin{equation}
 \varepsilon_{\rm pair,res}
 \le\min\left\{1,U\left(
       \varepsilon_{\rm MLDSA,EUF}^{\rm su}
       +\beta_{256}(q)
       +\delta_{\mathsf X}(q_*^{\rm res})
       +\delta_{\mathsf X}(q_*^{\rm ord})\right)\right\}.
 \label{eq:native-pair-ordinary-mldsa}
\end{equation}
The ML-DSA oracle ledger is a subset of the selected native pair owner's
ledger, hence of its actual-key ledger. No aggregation of distinct pair
owners with coincident ML-DSA components is assumed. Their original
generation and disclosure distribution is preserved. Cross-owner seed or
hedge-prefix coincidences are among the remaining shared-oracle queries
already counted in the selected-target comparison.

Equation~\eqref{eq:native-pair-ordinary-mldsa} supplies an ordinary-ML-DSA
instantiation of the signature term in
Proposition~\ref{prop:executed-source-attribution}, preserving its record,
KEM-correctness and joint-continuation accounting.
